
Discutere in questo capitolo come è stata progettata la soluzione al problema trattato nella tesi, indicando anche se sono stati valutati vari possibili approcci o soluzioni pre-esistenti e giustificando le proprie scelte. Descrivere quindi la soluzione vera e propria.

Nel caso sia stato sviluppato del software non triviale, è buona norma dedicargli tre sezioni:
\begin{itemize}
\item architettura dell'applicazione (interazioni con gli utenti e con altri sistemi, moduli logici, flussi dati interni ed esterni);
\item manuale dello sviluppatore (descrizione dei moduli, degli algoritmi, delle interfacce e delle strutture dati);
\item manuale utente (come installare ed usare il programma, interfacce, comandi, dati in input ed in output).
\end{itemize}
Nel caso di software molto voluminoso, queste tre sezioni possono diventare tre capitoli separati.

\chapter{Risultati}

Inserire in questo capitolo i risultati conseguiti, cercando di analizzarli -- se possibile -- in modo quantitativo.


\chapter{Conclusioni}

Qui si inseriscono brevi conclusioni sul lavoro svolto, senza ripetere inutilmente il sommario.

Si possono evidenziare i punti di forza e quelli di debolezza, nonché i possibili sviluppi futuri o attività da svolgere per migliorare i risultati.